\section{Exploratory Data Analysis}
Before training the classifier, we performed an exploratory data analysis (EDA) in order to gain a better
understanding of the feature distributions, potential outliers, and relationships between variables.
All visualizations were generated in Python using the \texttt{matplotlib} and \texttt{seaborn} libraries.

\subsection{Histograms of Numeric Features}

To gain an initial understanding of the numerical attributes in the dataset, histograms were generated for
the five key continuous features: \texttt{age}, \texttt{chol}, \texttt{oldpeak}, \texttt{thalach}, and \texttt{trestbps}.
These plots provide insights into the overall shape of each distribution, the presence of skewness, and the
frequency of extreme values. The following observations summarize the patterns apparent in the data.

\paragraph{Age.}
Figure~\ref{fig:hist-age} shows that the patients' ages range roughly from 30 to 80 years.
The distribution is close to a quasi-normal shape with a clear peak around 55--60 years.
Middle-aged and older adults constitute the majority of the dataset, while individuals younger than 40
appear far less frequently. This pattern is clinically expected, as the risk of heart disease
increases markedly with age. A mild decline in frequency is observed beyond age 65--70, reflecting the
lower proportion of very elderly patients in the dataset.

\begin{figure}[H]
	\centering
	\includegraphics[width=0.7\textwidth]{hist_age.png}
	\caption{Histogram of the \texttt{age} attribute.}
	\label{fig:hist-age}
\end{figure}

\paragraph{Cholesterol (\texttt{chol}).}
The histogram of \texttt{chol} (Figure~\ref{fig:hist-chol}) exhibits a moderately right-skewed distribution.
The majority of values lie between 200 and 260 mg/dl, a range commonly associated with borderline or high
cholesterol levels in clinical practice. A long upper tail extends to extreme values near 400--550 mg/dl,
reflecting the presence of a few patients with severe lipid abnormalities. These points represent medically
plausible outliers and are consistent with real-world cardiac risk factors.

\begin{figure}[H]
	\centering
	\includegraphics[width=0.7\textwidth]{hist_chol.png}
	\caption{Histogram of the \texttt{chol} attribute.}
	\label{fig:hist-chol}
\end{figure}

\paragraph{Oldpeak.}
The feature \texttt{oldpeak}, representing ST depression induced by exercise, has the most heavily
right-skewed distribution among all numeric attributes (Figure~\ref{fig:hist-oldpeak}). A large proportion
of patients have an \texttt{oldpeak} value close to zero, indicating minimal ST segment change.
The frequencies decrease rapidly as the value increases, with only a small number of observations above 3,
and rare extreme values reaching 5 or 6. This pattern aligns with clinical expectations, as significant
ST depression is relatively uncommon in the general cardiac patient population.

\begin{figure}[H]
	\centering
	\includegraphics[width=0.7\textwidth]{hist_oldpeak.png}
	\caption{Histogram of the \texttt{oldpeak} attribute.}
	\label{fig:hist-oldpeak}
\end{figure}

\paragraph{Thalach.}
The maximum heart rate (\texttt{thalach}) follows a nearly symmetric, bell-shaped distribution
(Figure~\ref{fig:hist-thalach}). Most patients achieve a heart rate between 150 and 170 beats per minute
during stress testing. Extremely low values (below 100) are rare, suggesting possible severe cardiac
conditions or incomplete exercise tests. This behavior of \texttt{thalach} is well aligned with clinical
observations across stress test populations.

\begin{figure}[H]
	\centering
	\includegraphics[width=0.7\textwidth]{hist_thalach.png}
	\caption{Histogram of the \texttt{thalach} attribute.}
	\label{fig:hist-thalach}
\end{figure}

\paragraph{Resting Blood Pressure (\texttt{trestbps}).}
As shown in Figure~\ref{fig:hist-trestbps}, \texttt{trestbps} displays a distribution slightly skewed
toward higher values. The central cluster of observations lies between 120 and 140 mmHg, corresponding
to mild or moderate hypertension, a common risk factor for heart disease. Higher values up to
approximately 180--200 mmHg appear infrequently but remain clinically realistic for untreated or
high-risk patients.

\begin{figure}[H]
	\centering
	\includegraphics[width=0.7\textwidth]{hist_trestbps.png}
	\caption{Histogram of the \texttt{trestbps} attribute.}
	\label{fig:hist-trestbps}
\end{figure}

Taken together, the histograms confirm that the numerical attributes exhibit a combination of
near-normal, moderately skewed, and heavily skewed distributions, each consistent with
real-world medical behavior. These observations further motivate the use of normalization
before applying distance-based algorithms such as KNN.

\subsection{Boxplots and Outliers}

Boxplots were generated for the five main numeric features in order to examine their spread, central
tendency, and potential outliers. These visualizations highlight the variability of each clinical measurement
and provide insights into whether extreme values are present. The following describes the observations
from each boxplot in detail.

\paragraph{Age.}
The boxplot of \texttt{age} (Figure~\ref{fig:boxplot-age}) shows that the majority of patients fall between
approximately 50 and 62 years of age, with a median around 56--58 years. No significant outliers are present,
indicating a relatively consistent age distribution. This pattern aligns well with clinical expectations,
as heart disease is most common among middle-aged and older adults.

\begin{figure}[H]
	\centering
	\includegraphics[width=0.7\textwidth]{boxplot_age.png}
	\caption{Boxplot of the \texttt{age} attribute.}
	\label{fig:boxplot-age}
\end{figure}

\paragraph{Cholesterol (\texttt{chol}).}
The boxplot of \texttt{chol} (Figure~\ref{fig:boxplot-chol}) reveals a number of high-value outliers,
with extreme cholesterol levels above 400 and even near 550 mg/dl. The interquartile range lies roughly
between 200 and 275 mg/dl, and the median is around 240--250 mg/dl. These outliers are medically plausible,
as severe hypercholesterolemia is a well-known risk factor for cardiac conditions.

\begin{figure}[H]
	\centering
	\includegraphics[width=0.7\textwidth]{boxplot_chol.png}
	\caption{Boxplot of the \texttt{chol} attribute.}
	\label{fig:boxplot-chol}
\end{figure}

\paragraph{Oldpeak.}
The boxplot of \texttt{oldpeak} (Figure~\ref{fig:boxplot-oldpeak}) shows a strongly right-skewed
distribution, with the majority of values located between 0 and 1.5. A few substantial outliers appear above 5,
indicating rare but clinically significant cases of severe ST depression during exercise testing. These extreme
values contribute to the heavy-tailed structure observed in the histogram of this feature.

\begin{figure}[H]
	\centering
	\includegraphics[width=0.7\textwidth]{boxplot_oldpeak.png}
	\caption{Boxplot of the \texttt{oldpeak} attribute.}
	\label{fig:boxplot-oldpeak}
\end{figure}

\paragraph{Thalach.}
In the boxplot of \texttt{thalach} (Figure~\ref{fig:boxplot-thalach}), the central mass of the data lies between
approximately 140 and 170 beats per minute, with a median near 155 bpm. A single low outlier around 80 bpm is
visible, which may correspond to severe cardiac limitations or an incomplete stress test. Upper values extend to
about 200 bpm, which is physiologically plausible in stress testing scenarios.

\begin{figure}[H]
	\centering
	\includegraphics[width=0.7\textwidth]{boxplot_thalach.png}
	\caption{Boxplot of the \texttt{thalach} attribute.}
	\label{fig:boxplot-thalach}
\end{figure}

\paragraph{Resting Blood Pressure (\texttt{trestbps}).}
The boxplot of \texttt{trestbps} (Figure~\ref{fig:boxplot-trestbps}) shows a typical clinical distribution,
with most values in the 120--140 mmHg range and a median around 130--135 mmHg. Several pronounced outliers
appear at 170--200 mmHg, indicating cases of severe hypertension. These observations are consistent with known
risk patterns in cardiac patients.

\begin{figure}[H]
	\centering
	\includegraphics[width=0.7\textwidth]{boxplot_trestbps.png}
	\caption{Boxplot of the \texttt{trestbps} attribute.}
	\label{fig:boxplot-trestbps}
\end{figure}

Overall, the boxplots confirm the presence of meaningful medical outliers in features such as
\texttt{chol}, \texttt{oldpeak}, and \texttt{trestbps}, while attributes like \texttt{age} and
\texttt{thalach} exhibit relatively stable distributions. Since KNN is generally robust to a small
number of well-separated outliers and normalization was applied during preprocessing, these extreme
values were retained rather than removed.

\subsection{Q--Q Plots}

To examine how closely the numeric features follow a normal distribution, Q--Q plots were generated for
three representative attributes: \texttt{age}, \texttt{chol}, and \texttt{thalach}. In each plot, the
sorted sample values are compared against the theoretical quantiles of a normal distribution with the
same mean and standard deviation. Deviations from the diagonal line indicate departures from normality.

\paragraph{Age.}
The Q--Q plot for \texttt{age} (Figure~\ref{fig:qq-age}) shows that most points lie close to the diagonal
line, indicating that the central portion of the distribution is approximately normal. However, noticeable
deviations appear in both tails: younger ages fall slightly below the reference line, while older ages
(above roughly 75 years) lie somewhat above it. This pattern suggests a mildly heavy-tailed distribution
but overall a fairly symmetric structure. Such behavior is typical for age in clinical datasets.

\begin{figure}[H]
	\centering
	\includegraphics[width=0.7\textwidth]{qq_age.png}
	\caption{Q--Q plot of the \texttt{age} attribute.}
	\label{fig:qq-age}
\end{figure}

\paragraph{Cholesterol (\texttt{chol}).}
The Q--Q plot of \texttt{chol} (Figure~\ref{fig:qq-chol}) reveals much stronger deviations, particularly
in the upper tail. While the central region aligns reasonably well with the diagonal, several extreme
values above 400 mg/dl rise sharply above the theoretical normal quantiles. These points correspond to
clinically plausible cases of severe hypercholesterolemia. The resulting right-skewed behavior confirms
that \texttt{chol} does not follow a normal distribution and contains meaningful outliers.

\begin{figure}[H]
	\centering
	\includegraphics[width=0.7\textwidth]{qq_chol.png}
	\caption{Q--Q plot of the \texttt{chol} attribute.}
	\label{fig:qq-chol}
\end{figure}

\paragraph{Thalach.}
The Q--Q plot of \texttt{thalach} (Figure~\ref{fig:qq-thalach}) shows that the central portion of the
distribution is approximately linear, indicating an almost normal pattern. However, the lower tail contains
a few markedly low values (around 70--90 bpm), which fall well below the diagonal. Similarly, some values
above 200 bpm appear slightly above the theoretical line. These deviations suggest mild skewness driven by
rare low and high measurements, consistent with clinical variations observed in exercise testing.

\begin{figure}[H]
	\centering
	\includegraphics[width=0.7\textwidth]{qq_thalach.png}
	\caption{Q--Q plot of the \texttt{thalach} attribute.}
	\label{fig:qq-thalach}
\end{figure}

Overall, the Q--Q plots confirm that none of the examined features follow a perfectly normal distribution.
Each attribute exhibits tail deviations or skewness, with \texttt{chol} showing the strongest departure due
to clinically meaningful high outliers. This is not problematic for the KNN classifier, which does not rely
on normality assumptions.

\subsection{Correlation Matrix}

To investigate linear relationships among the numeric features and their association with the target variable,
the Pearson correlation matrix was computed and visualized as a heatmap (Figure~\ref{fig:corr-matrix}).
The matrix reveals several clinically meaningful patterns and highlights which attributes are most relevant
for distinguishing between patients with and without heart disease.

\begin{figure}[H]
	\centering
	\includegraphics[width=0.75\textwidth]{corr_matrix.png}
	\caption{Correlation matrix of selected numeric features and the target variable.}
	\label{fig:corr-matrix}
\end{figure}

The strongest correlations with the target variable are observed for \texttt{thalach} and \texttt{oldpeak}.
The feature \texttt{thalach} (maximum heart rate achieved) shows a moderately positive correlation ($r = 0.42$),
indicating that individuals who achieve higher heart rates during exercise are less likely to have heart disease.
Conversely, \texttt{oldpeak} exhibits a moderately negative correlation with the target ($r = -0.44$), consistent
with the clinical understanding that larger ST depression values are strongly associated with myocardial ischemia
and a higher likelihood of disease. Age also shows a mild negative correlation with the target ($r = -0.22$),
reflecting the increased prevalence of heart disease among older patients, although the relationship is not strictly linear.

Other features such as \texttt{trestbps} and \texttt{chol} display only weak correlations with the target,
suggesting that—within this dataset—they contribute less predictive power on their own but may still play a role
when combined with other variables in a multivariate context.

Several inter-feature relationships are also noteworthy. Age and \texttt{thalach} are moderately negatively correlated
($r = -0.38$), meaning that older individuals typically achieve lower maximum heart rates during stress testing.
Likewise, \texttt{oldpeak} and \texttt{thalach} show a negative correlation ($r = -0.35$), reflecting the reduced
exercise tolerance and performance typically seen in patients with ischemic responses.

Overall, the correlation matrix confirms that heart disease in this dataset is most strongly associated with
exercise-related features such as \texttt{thalach} and \texttt{oldpeak}, while resting measurements such as
\texttt{trestbps} and \texttt{chol} exhibit weaker linear dependencies with the target. These insights also highlight
why distance-based models like KNN, which account for interactions among multiple features simultaneously,
can still perform effectively despite the limited strength of individual linear correlations.

\subsection{Pair Plots (Scatterplot Matrix)}

To investigate multivariate relationships among the key numeric attributes, a pair plot was generated for
\texttt{age}, \texttt{chol}, \texttt{thalach}, and \texttt{oldpeak}, with points colored according to the
binary target variable. This visualization provides a simultaneous view of the marginal distributions of
each feature as well as their pairwise interactions (Figure~\ref{fig:pairplot-selected}).

\begin{figure}[H]
	\centering
	\includegraphics[width=0.85\textwidth]{pairplot_selected.png}
	\caption{Pair plot for selected numeric features, colored by the target class.}
	\label{fig:pairplot-selected}
\end{figure}

The diagonal histograms show that \texttt{age} and \texttt{thalach} follow approximately symmetric
distributions, whereas \texttt{chol} and \texttt{oldpeak} are noticeably right-skewed. In terms of class
distribution, patients with heart disease (target = 1) appear slightly more often in higher age ranges,
but the effect is modest. In contrast, clear class differences emerge in the distributions of
\texttt{thalach} and \texttt{oldpeak}: individuals without heart disease typically achieve higher maximum
heart rates, while those with heart disease tend to exhibit larger ST depression values.

The pairwise scatter plots reveal several important patterns. The relationship between \texttt{age} and
\texttt{thalach} displays a clear negative trend reflecting the expected physiological decline in exercise
capacity with age. Moreover, the combination of \texttt{thalach} and \texttt{oldpeak} provides the strongest
visual separation between the two classes: non-diseased individuals cluster around high \texttt{thalach}
and near-zero \texttt{oldpeak}, whereas diseased individuals tend to show lower \texttt{thalach} values and
elevated \texttt{oldpeak}. In contrast, \texttt{chol} exhibits substantial overlap between the two classes
in all pairwise projections, indicating that it is a weaker standalone discriminating feature.

Overall, the pair plots confirm that heart disease in this dataset is most strongly associated with exercise
performance indicators (\texttt{thalach}, \texttt{oldpeak}), while features such as \texttt{chol} and
\texttt{age} show weaker or less consistent separation. These multivariate patterns justify the use of a
distance-based classifier such as KNN, which can exploit such relationships even when they are nonlinear or
not clearly separable along any single dimension.

\subsection{Count Plots of Categorical Features}

To better understand the distribution of categorical variables in the dataset, count plots were generated for
\texttt{cp} (chest pain type), \texttt{exang} (exercise-induced angina), \texttt{sex}, and the binary
\texttt{target} variable. These plots illustrate the frequency of each category and help identify any potential
imbalances in the dataset.

\paragraph{Chest Pain Type (\texttt{cp}).}
Figure~\ref{fig:count-cp} shows that chest pain type is far from uniformly distributed. The category
\texttt{cp} = 0 (typical angina) is the most common with nearly 500 observations, while \texttt{cp} = 3
(asymptomatic) is the least frequent with fewer than 100 samples. The other two categories fall between these
extremes. This distribution reflects real-world clinical observations, where typical angina is often the primary
presenting symptom among cardiac patients.

\begin{figure}[H]
	\centering
	\includegraphics[width=0.7\textwidth]{count_cp.png}
	\caption{Count plot of the \texttt{cp} (chest pain type) attribute.}
	\label{fig:count-cp}
\end{figure}

\paragraph{Exercise-Induced Angina (\texttt{exang}).}
As shown in Figure~\ref{fig:count-exang}, the majority of patients (\texttt{exang} = 0) did not experience
exercise-induced angina during stress testing. A smaller subset of roughly one-third of the dataset consists
of patients who did experience angina (\texttt{exang} = 1). This imbalance is expected, as not all individuals
display angina symptoms even when underlying disease is present.

\begin{figure}[H]
	\centering
	\includegraphics[width=0.7\textwidth]{count_exang.png}
	\caption{Count plot of the \texttt{exang} attribute.}
	\label{fig:count-exang}
\end{figure}

\paragraph{Sex.}
Figure~\ref{fig:count-sex} demonstrates that the dataset contains a higher proportion of male patients
(\texttt{sex} = 1) compared to female patients (\texttt{sex} = 0), with an approximate ratio of 2:1.
This distribution reflects the higher prevalence of heart disease among men in many clinical populations.

\begin{figure}[H]
	\centering
	\includegraphics[width=0.7\textwidth]{count_sex.png}
	\caption{Count plot of the \texttt{sex} attribute.}
	\label{fig:count-sex}
\end{figure}

\paragraph{Target Variable.}
The class distribution shown in Figure~\ref{fig:count-target} indicates that the dataset is nearly balanced,
with similar numbers of samples for \texttt{target} = 0 (no heart disease) and \texttt{target} = 1
(presence of heart disease). This balance is advantageous for training classification models, as it reduces
the risk of bias toward either class and removes the need for rebalancing techniques such as oversampling or
class weighting.

\begin{figure}[H]
	\centering
	\includegraphics[width=0.7\textwidth]{count_target.png}
	\caption{Count plot of the \texttt{target} variable.}
	\label{fig:count-target}
\end{figure}

Overall, the count plots reveal realistic clinical distributions across the categorical features. Although some
categories—such as chest pain type and sex—exhibit notable imbalances, the balanced target distribution ensures
that classification models such as KNN can be trained effectively without requiring adjustments for class imbalance.